\scanpage{122}
\mspage{113}
% archive sheet number: 122

\deleted{2°) Dans $\delta\Psi=\partial\Psi$,}

Les sommets \deleted{isolés sont} internes sont \deleted{ceux de}\inserted{les} éléments de
\[
\Delta_{\mathrm{int}}=(\Delta\deleted{\unsure{\epsilon}}\setminus\partial\Psi]\cap\mathrm{Omb}(\Psi)=\Delta\cap\mathrm{Omb}^{\circ}(\Psi)
\]
Les sommets isolés sont \deleted{ceux de} les autres sommets $\deleted{\Delta\setminus\partial\Psi}$ de $\Delta\setminus\partial\Psi$,
\[
\Delta_{\mathrm{is}}=\deleted{\unsure{\Delta\setminus(\partial\Psi\cup\Delta}}(\Delta\setminus\partial\Psi)\setminus\Delta_{\mathrm{int}}
\]
\deleted{2°) \unsure{Distinguer}, \unsure{donc donner}, dans $\Delta_{\mathrm{int}}$,}
\deleted{une partie}
donc on a une décomposition
\[
\Delta=\partial\Psi\amalg\Delta_{\mathrm{int}}\amalg\Delta_{\mathrm{is}}
\]

2°) Se donner, dans $\Delta_{\mathrm{int}}$, les sommets qui appartiennent à $\Phi_0$ (sommets internes propres, ou sommets de \unsure{redondance} ou de \unsure{morcèlement|morcellement}), ou ce qui revient au $=$, les autres : \unsure{les} sommets lacunaires
\[
\boxed{\Delta_{\mathrm{lac}}}\subset\Delta_{\mathrm{int}}
\]
de sorte que
\[
\Delta_{\mathrm{int}}=\Delta_{\mathrm{lac}}\amalg\Delta_{\mathrm{red}}\,.
\]

3°) Dans $\partial\Psi$\inserted{$(=\partial\Phi)$}, dire quels sommets sont dans $\Phi_0$ \deleted{et lesquels} (sommets bords propres de $\Phi$) et lesquels non, i.e. se donner une partie
\[
\begin{aligned}
&\deleted{\unsure{\partial_0\Psi|\delta_0\Psi}}\\
&\unsure{\partial_0|\delta_0}\subset\partial\Psi\,.
\end{aligned}
\]

En résumé :
